\documentclass[11pt]{article}
\usepackage{ctex}
\usepackage{amsmath,amssymb,mathtools}
\usepackage{tikz}
\usepackage{tikz-cd}
\usetikzlibrary{positioning,arrows.meta,calc,fit}
\usepackage{geometry}
\usepackage{booktabs}
\usepackage{xcolor}
\geometry{a3paper}

% ============================================================
% 常用宏
% ============================================================
\newcommand{\RR}{\mathbb R}
\newcommand{\ZZ}{\mathbb Z}
\newcommand{\Rbig}{\mathcal R_N}
\newcommand{\Rsmall}{\mathcal R_n}
\newcommand{\Sn}{\mathcal S_n}
\newcommand{\Snh}{\mathcal S_{n/2}}
\newcommand{\Pe}{P_{\mathrm{even}}}
\newcommand{\Po}{P_{\mathrm{odd}}}
\newcommand{\ct}{\mathbf{ct}}
\newcommand{\ctP}{\mathbf{ct}_P}
\newcommand{\cte}{\mathbf{ct}_{\mathrm{even}}}
\newcommand{\cto}{\mathbf{ct}_{\mathrm{odd}}}
\newcommand{\ctoy}{\mathbf{ct}_{\mathrm{odd},Y}}
\newcommand{\KSK}{\mathsf{KSK}}
\newcommand{\KS}{\mathsf{KeySwitch}}
\newcommand{\Comp}{\mathsf{Compress}}
\newcommand{\Split}{\mathsf{Split}}
\newcommand{\Merge}{\mathsf{Merge}}
\newcommand{\Auto}{\mathsf{Auto}}
\newcommand{\MulPt}{\mathsf{MulPt}}
\newcommand{\Add}{\mathsf{Add}}
\newcommand{\Sub}{\mathsf{Sub}}
\newcommand{\Enc}{\mathsf{Enc}}
\newcommand{\Dec}{\mathsf{Dec}}
\newcommand{\Level}{\mathsf{Level}}
\newcommand{\ksnoise}{\eta_{\mathrm{ks}}}
\newcommand{\rep}{\mathrel{\dashv\vdash}}

\title{CKKS 中的同态 FFT 奇偶拆分与大环到小环压缩：图解笔记}
\author{}
\date{}

\begin{document}
\maketitle

\tableofcontents
\newpage


\section*{阅读指南}

本文把两个推导放在同一套图景里理解。

\begin{itemize}
  \item \textbf{文档 1} 解决的是：如何在密文状态下对多项式做 FFT 式奇偶拆分。
  \item \textbf{文档 2} 解决的是：当拆出来的结果形如 $f(X^2)$ 时，如何把它从大环密文压缩成小环密文 $f(Y)$。
\end{itemize}

一句话概括二者关系：
\[
\boxed{\text{文档 1：拆出 } f(X^2)}
\qquad
\boxed{\text{文档 2：把 } f(X^2) \text{ 压成 } f(Y)}.
\]

\section{总览：两个文档如何连接}

先看整体流程。文档 1 先把一个大环密文拆成偶分支和奇分支；文档 2 再把每个形如 $f(X^2)$ 的分支压到小环。

\[
\begin{tikzcd}[column sep=large,row sep=large]
\text{大环密文 } \ctP
  \arrow[r,"\text{文档 1：同态奇偶拆分}"]
&
\begin{matrix}
\cte \text{ encrypts } \Pe(X^2) \\
\cto \text{ encrypts } \Po(X^2)
\end{matrix}
  \arrow[r,"\text{文档 2：大环到小环压缩}"]
&
\begin{matrix}
\ct'_e \text{ encrypts } \Pe(Y) \\
\ct'_o \text{ encrypts } \Po(Y)
\end{matrix}
\end{tikzcd}
\]

这里的核心点是：文档 1 只让明文语义变成半长度结构，结果仍然存放在大环里；文档 2 则进一步把密文结构也降到小环。

\newpage


\section{符号说明}

本节集中列出本文使用的主要符号。整体上，本文涉及两类操作：
一类是文档 1 中的同态 FFT 式奇偶拆分；
另一类是文档 2 中的大环密文到小环密文压缩。

\subsection{环与变量}

\begin{center}
\begin{tabular}{lll}
\toprule
符号 & 含义 & 说明 \\
\midrule
\(X\) & 大环变量 & CKKS 大环中的变量 \\
\(Y\) & 小环变量 & 压缩后小环中的变量 \\
\(T\) & 分支多项式变量 & 用于表示偶分支、奇分支的抽象变量 \\
\(N\) & 大环维度 & 通常为 \(2\) 的幂 \\
\(n\) & 半长度参数 & 通常满足 \(n\mid N\) \\
\(k\) & 嵌入步长 & \(k=\frac{N}{2n}\) \\
\(\Rbig\) & 大环 & \(\mathcal R_N=\mathbb R[X]/(X^N+1)\) 或 \(\mathbb Z_q[X]/(X^N+1)\) \\
\(\Rsmall\) & 小环 & \(\mathcal R_n=\mathbb R[Y]/(Y^n+1)\) 或 \(\mathbb Z_q[Y]/(Y^n+1)\) \\
\(\Sn\) & 文档 1 中的源环 & \(\mathcal S_n=\mathbb R[Y]/(Y^{2n}+1)\) \\
\(\Snh\) & 文档 1 中的半长度环 & \(\mathcal S_{n/2}=\mathbb R[T]/(T^n+1)\) \\
\bottomrule
\end{tabular}
\end{center}

\subsection{多项式与奇偶分支}

\begin{center}
\begin{tabular}{lll}
\toprule
符号 & 含义 & 说明 \\
\midrule
\(P(Y)\) & 原始多项式 & \(P(Y)=\sum_{j=0}^{2n-1}p_jY^j\) \\
\(\Pe(T)\) & 偶分支多项式 & \(\Pe(T)=\sum_{i=0}^{n-1}p_{2i}T^i\) \\
\(\Po(T)\) & 奇分支多项式 & \(\Po(T)=\sum_{i=0}^{n-1}p_{2i+1}T^i\) \\
\(\Pe(Y^2)\) & 偶部嵌回原变量 & 由 \(P(Y)\) 的偶数次系数组成 \\
\(Y\Po(Y^2)\) & 带 \(Y\) 因子的奇部 & 由 \(P(Y)\) 的奇数次系数组成 \\
\(\Pe(X^{2k})\) & 大环中的偶分支 & 将 \(Y=X^k\) 后得到 \\
\(\Po(X^{2k})\) & 大环中的归一化奇分支 & 去掉前面的 \(X^k\) 因子后得到 \\
\(X^k\Po(X^{2k})\) & 大环中的带因子奇分支 & 尚未去掉 \(X^k\) 的奇分支 \\
\bottomrule
\end{tabular}
\end{center}

核心恒等式为
\[
P(Y)=\Pe(Y^2)+Y\Po(Y^2),
\]
\[
P(-Y)=\Pe(Y^2)-Y\Po(Y^2).
\]
因此
\[
\Pe(Y^2)=\frac{P(Y)+P(-Y)}{2},
\qquad
Y\Po(Y^2)=\frac{P(Y)-P(-Y)}{2}.
\]

\subsection{嵌入、自同构与公开映射}

\begin{center}
\begin{tabular}{lll}
\toprule
符号 & 含义 & 说明 \\
\midrule
\(\eta_k\) & 嵌入映射 & \(\eta_k(P(Y))=P(X^k)\) \\
\(\eta_{2k}\) & 半长度分支嵌入 & \(A(T)\mapsto A(X^{2k})\) \\
\(\iota_n\) & 取负变量映射 & \(\iota_n(P(Y))=P(-Y)\) \\
\(\sigma_a\) & 大环自同构 & \(\sigma_a(F(X))=F(X^a)\)，其中 \(a\) 为奇数 \\
\(\sigma_{1+2n}\) & 实现 \(X^k\mapsto -X^k\) 的自同构 & 因为 \(\sigma_{1+2n}(X^k)=X^{k+N}=-X^k\) \\
\(\tau\) & 偶次系数提取映射 & \(\tau(F)(Y)=\sum_{i=0}^{n-1}f_{2i}Y^i\) \\
\(X^{-1}\) & 大环中 \(X\) 的逆元 & 当 \(N=2n\) 时，\(X^{-1}=-X^{2n-1}\) \\
\bottomrule
\end{tabular}
\end{center}

其中，若
\[
F(X)=\sum_{j=0}^{2n-1}f_jX^j,
\]
则
\[
\tau(F)(Y)=\sum_{i=0}^{n-1}f_{2i}Y^i.
\]
即 \(\tau\) 是“取偶数下标系数并重新编号”的公开线性映射。

\subsection{CKKS 密文与基本操作}

\begin{center}
\begin{tabular}{lll}
\toprule
符号 & 含义 & 说明 \\
\midrule
\(\ct\) & 一般密文 & 通常写作 \(\ct=(c_0,c_1)\) \\
\(\ctP\) & 加密 \(P(X^k)\) 的密文 & 文档 1 的输入密文 \\
\(\ct_-\) & 加密 \(P(-X^k)\) 的密文 & \(\ct_-=\Auto_{1+2n}(\ctP)\) \\
\(\cte\) & 偶分支密文 & 加密 \(\Pe(X^{2k})\) \\
\(\ctoy\) & 带因子奇分支密文 & 加密 \(X^k\Po(X^{2k})\) \\
\(\cto\) & 归一化奇分支密文 & 加密 \(\Po(X^{2k})\) \\
\(\Auto_a\) & 密文自同构操作 & 对应明文层 \(X\mapsto X^a\) \\
\(\Add\) & 密文加法 & 明文语义为相加 \\
\(\Sub\) & 密文减法 & 明文语义为相减 \\
\(\MulPt{u}\) & 明文--密文乘法 & 将密文消息乘以公开明文 \(u\) \\
\(\rep\) & 密文语义记号 & \(\ct\rep F\) 表示密文 \(\ct\) 解密后语义上为 \(F\) \\
\bottomrule
\end{tabular}
\end{center}

文档 1 中常用操作为
\[
\ct_-:=\Auto_{1+2n}(\ctP),
\]
\[
\cte:=\frac{1}{2}(\ctP+\ct_-),
\]
\[
\ctoy:=\frac{1}{2}(\ctP-\ct_-),
\]
\[
\cto:=(-X^{N-k})\ctoy.
\]

对应语义为
\[
\cte\rep \Pe(X^{2k}),
\qquad
\ctoy\rep X^k\Po(X^{2k}),
\qquad
\cto\rep \Po(X^{2k}).
\]

\subsection{文档 2 中的大环到小环压缩符号}

\begin{center}
\begin{tabular}{lll}
\toprule
符号 & 含义 & 说明 \\
\midrule
\(c_0(X),c_1(X)\) & 大环密文分量 & 输入密文 \(\ct=(c_0,c_1)\) 的两个公开多项式 \\
\(s(X)\) & 大环私钥 & 解密时满足 \(c_0+s c_1\approx \Delta m(X^2)\) \\
\(m(Y)\) & 小环消息多项式 & 目标消息 \\
\(m(X^2)\) & 大环中的嵌入消息 & 只含 \(X\) 的偶次幂 \\
\(\Delta\) & CKKS 缩放因子 & 明文缩放比例 \\
\(e(X)\) & 原始大环噪声 & 输入密文的解密误差 \\
\(A\) & \(c_0\) 的偶部 & \(A=\tau(c_0)\) \\
\(B\) & \(c_1\) 的偶部 & \(B=\tau(c_1)\) \\
\(C\) & \(c_0\) 的奇部重编号 & \(C=\tau(X^{-1}c_0)\) \\
\(D\) & \(c_1\) 的奇部重编号 & \(D=\tau(X^{-1}c_1)\) \\
\(s_0(Y)\) & 私钥偶部 & \(s(X)=s_0(X^2)+Xs_1(X^2)\) \\
\(s_1(Y)\) & 私钥奇部 & 同上 \\
\(s'(Y)\) & 目标小环私钥 & key switching 后绑定的普通小环私钥 \\
\(\nu_0\) & 偶部噪声 & 来自 \(e_{\mathrm{even}}\) \\
\(\nu_1\) & 奇部噪声 & 来自 \(e_{\mathrm{odd}}\) \\
\(\eta_{\mathrm{ks}}\) & key switching 噪声 & 密钥切换新增误差 \\
\bottomrule
\end{tabular}
\end{center}

文档 2 中有
\[
c_0(X)=A(X^2)+XC(X^2),
\]
\[
c_1(X)=B(X^2)+XD(X^2),
\]
\[
s(X)=s_0(X^2)+Xs_1(X^2).
\]

代入解密关系
\[
c_0+s c_1=\Delta m(X^2)+e(X)
\]
后，得到小环中的两个方程：
\[
A+s_0B+Y s_1D=\Delta m(Y)+\nu_0,
\]
\[
C+s_0D+s_1B=\nu_1.
\]

第一行携带目标消息，第二行是近似零关系。

\subsection{模块密文与 key switching 符号}

\begin{center}
\begin{tabular}{lll}
\toprule
符号 & 含义 & 说明 \\
\midrule
\((t,u,v)\) & 三元素模块密文 & 满足 \(t+s_0u+s_1v\approx \mathrm{msg}\) \\
\((A,B,YD)\) & 本文主模块密文 & 满足 \(A+s_0B+s_1(YD)\approx \Delta m(Y)\) \\
\(G\) & gadget 基 & 用于 gadget 分解 \\
\(L\) & gadget 分解长度 & 分解层数 \\
\(u_\ell,v_\ell\) & gadget digits & \(u=\sum_\ell u_\ell G^\ell,\ v=\sum_\ell v_\ell G^\ell\) \\
\(a_{0,\ell},a_{1,\ell}\) & key switching 随机项 & 公开切换密钥的一部分 \\
\(b_{0,\ell},b_{1,\ell}\) & key switching 加密项 & 公开切换密钥的一部分 \\
\(\epsilon_{0,\ell},\epsilon_{1,\ell}\) & 切换密钥噪声 & key switching key 中的小误差 \\
\(d_0,d_1\) & 输出小环密文分量 & 满足 \(d_0+s'd_1\approx \mathrm{msg}\) \\
\(\KS\) & key switching 操作 & 将私钥从 \((s_0,s_1)\) 切到 \(s'\) \\
\bottomrule
\end{tabular}
\end{center}

本文采用解密约定
\[
d_0+s'd_1.
\]
因此正确的切换密钥形式为
\[
b_{0,\ell}=G^\ell s_0-a_{0,\ell}s'+\epsilon_{0,\ell},
\]
\[
b_{1,\ell}=G^\ell s_1-a_{1,\ell}s'+\epsilon_{1,\ell}.
\]

若
\[
u=\sum_{\ell=0}^{L-1}u_\ell G^\ell,
\qquad
v=\sum_{\ell=0}^{L-1}v_\ell G^\ell,
\]
则 key switching 后的新增噪声为
\[
\eta_{\mathrm{ks}}
=
\sum_{\ell=0}^{L-1}u_\ell\epsilon_{0,\ell}
+
\sum_{\ell=0}^{L-1}v_\ell\epsilon_{1,\ell}.
\]

对于文档 2 的主输出，
\[
(t,u,v)=(A,B,YD),
\]
因此最终小环密文满足
\[
d_0+s'd_1
=
\Delta m(Y)+\nu_0+\eta_{\mathrm{ks}}.
\]

\subsection{噪声与层数符号}

\begin{center}
\begin{tabular}{lll}
\toprule
符号 & 含义 & 说明 \\
\midrule
\(e(X)\) & 输入大环噪声 & 可拆为偶部和奇部 \\
\(e_{\mathrm{even}}(X^2)\) & 原噪声偶部 & 会进入输出噪声 \\
\(Xe_{\mathrm{odd}}(X^2)\) & 原噪声奇部 & 不进入主输出第一行 \\
\(\nu_0\) & 小环偶部噪声 & \(\nu_0=\tau(e_{\mathrm{even}})\) \\
\(\nu_1\) & 小环奇部噪声 | 对应近似零关系 \\
\(\eta_{\mathrm{ks}}\) & key switching 噪声 & 压缩时新增 \\
\(e_{\mathrm{out}}\) & 最终输出噪声 & \(e_{\mathrm{out}}=\nu_0+\eta_{\mathrm{ks}}\) \\
\(\ell\) & 模数链层数索引 & level \(\ell\) \\
\(q_\ell\) & 当前模数 & 当前 level 下的密文模数 \\
\bottomrule
\end{tabular}
\end{center}

噪声流向可概括为
\[
e(X)=e_{\mathrm{even}}(X^2)+Xe_{\mathrm{odd}}(X^2),
\]
\[
e_{\mathrm{out}}=\nu_0+\eta_{\mathrm{ks}}.
\]

也就是说，大环到小环压缩保留原噪声的偶部，丢掉不承载目标消息的奇部噪声，同时新增 key switching 噪声。

\subsection{常用操作的层数消耗}

\begin{center}
\begin{tabular}{llll}
\toprule
操作 & 作用 & 默认层数消耗 & 说明 \\
\midrule
\(\Auto_a\) & 自同构 / rotation & 0 & 增噪声，需要 automorphism key \\
\(\Add,\Sub\) & 加法 / 减法 & 0 & 要求输入在同一 level \\
乘 \(1/2\) & 精确常数乘法 & 0 & 若作为 scale 1 plaintext 处理 \\
乘 \(X^k\) 或 \(X^{-1}\) & 单项式移位 & 0 & 循环移位加符号变化 \\
\(\tau\) & 公开系数提取 & 0 & 不属于同态运算 \\
gadget decomposition & 系数分解 & 0 & 不改变密文 level \\
key switching & 切换绑定私钥 & 0 & 通常不掉层，但增加噪声 \\
rescale / mod-down & 降模或缩放调整 & 1 或更多 & 主动降几层就消耗几层 \\
密文--密文乘法 & 同态乘法 & 通常 1 & 本文主流程中没有该操作 \\
\bottomrule
\end{tabular}
\end{center}

\newpage

\section{文档 1：同态 FFT 奇偶拆分}

\subsection{明文层的奇偶拆分}

设
\[
P(Y)=\sum_{j=0}^{2n-1}p_jY^j.
\]
定义
\[
\Pe(T)=\sum_{i=0}^{n-1}p_{2i}T^i,
\qquad
\Po(T)=\sum_{i=0}^{n-1}p_{2i+1}T^i.
\]
那么有精确恒等式
\[
P(Y)=\Pe(Y^2)+Y\Po(Y^2),
\]
以及
\[
P(-Y)=\Pe(Y^2)-Y\Po(Y^2).
\]
因此
\[
\Pe(Y^2)=\frac{P(Y)+P(-Y)}{2},
\qquad
Y\Po(Y^2)=\frac{P(Y)-P(-Y)}{2}.
\]

下面这张图表达了：只要能得到 $P(-Y)$，就能通过半和、半差得到偶分支和奇分支。

\[
\begin{tikzcd}[column sep=huge,row sep=large]
P(Y)
  \arrow[r,"Y\mapsto -Y"]
  \arrow[d,"\frac{1}{2}(P(Y)+P(-Y))"']
  \arrow[dd,"\frac{1}{2}(P(Y)-P(-Y))"', bend right=55]
&
P(-Y)
\\
\Pe(Y^2)
\\
Y\Po(Y^2)
\end{tikzcd}
\]

\subsection{嵌入到 CKKS 大环中}

令
\[
\Sn=\frac{\RR[Y]}{(Y^{2n}+1)},
\qquad
\Rbig=\frac{\RR[X]}{(X^N+1)},
\qquad
k=\frac{N}{2n}.
\]
通过
\[
Y=X^k
\]
把 $\Sn$ 嵌入到 $\Rbig$。因为
\[
(X^k)^{2n}=X^N=-1,
\]
所以这个嵌入是良定义的。

文档 1 的关键是用大环自同构实现
\[
X^k\mapsto -X^k.
\]
取
\[
\sigma_{1+2n}:X\mapsto X^{1+2n}.
\]
则
\[
\sigma_{1+2n}(X^k)=X^{k(1+2n)}=X^{k+N}=-X^k.
\]

图示如下：

\[
\begin{tikzcd}[column sep=huge,row sep=large]
\ctP
  \arrow[r,"\Auto_{1+2n}"]
  \arrow[d,dashed,"\text{decrypts to}"']
&
\ct_{-}
  \arrow[d,dashed,"\text{decrypts to}"]
\\
P(X^k)
  \arrow[r,"X^k\mapsto -X^k"']
&
P(-X^k)
\end{tikzcd}
\]

\subsection{密文层正向拆分}

输入密文满足
\[
\ctP \rep P(X^k).
\]
先做自同构：
\[
\ct_-:=\Auto_{1+2n}(\ctP),
\qquad
\ct_-\rep P(-X^k).
\]
然后计算
\[
\cte:=\frac{1}{2}(\ctP+\ct_-),
\]
\[
\ctoy:=\frac{1}{2}(\ctP-\ct_-).
\]
于是
\[
\cte\rep \Pe(X^{2k}),
\qquad
\ctoy\rep X^k\Po(X^{2k}).
\]
如果希望去掉奇分支前面的 $X^k$，继续乘
\[
(X^k)^{-1}=-X^{N-k},
\]
得到
\[
\cto:=(-X^{N-k})\ctoy,
\qquad
\cto\rep \Po(X^{2k}).
\]

流程图如下：

```latex
\[
\begin{tikzcd}[column sep=large,row sep=large]
& \ctP
  \arrow[dl, "\mathrm{id}"']
  \arrow[dr, "\Auto_{1+2n}"]
&
\\
\ctP
  \arrow[dr, "+"']
  \arrow[ddr, "-"', bend right=20]
&
&
\ct_{-}
  \arrow[dl, "+"]
  \arrow[ddl, "-", bend left=20]
\\
&
\frac{1}{2}(\ctP+\ct_{-})
  \arrow[d, dashed, "\Pe(X^{2k})"]
&
\\
&
\frac{1}{2}(\ctP-\ct_{-})
  \arrow[d, dashed, "X^k\Po(X^{2k})"]
&
\\
&
(-X^{N-k})\cdot \frac{1}{2}(\ctP-\ct_{-})
  \arrow[d, dashed, "\Po(X^{2k})"]
&
\\
& {} &
\end{tikzcd}
\]
```

\subsection{密文层重构}

如果没有做小环压缩，重构非常直接。

若已有
\[
\cte\rep \Pe(X^{2k}),
\qquad
\cto\rep \Po(X^{2k}),
\]
则
\[
\Merge(\cte,\cto)=\cte+X^k\cto,
\]
满足
\[
\Merge(\cte,\cto)\rep \Pe(X^{2k})+X^k\Po(X^{2k})=P(X^k).
\]

若奇分支保留了 $X^k$ 因子，即
\[
\ctoy\rep X^k\Po(X^{2k}),
\]
则直接相加：
\[
\Merge(\cte,\ctoy)=\cte+\ctoy.
\]

图示如下：

\[
\begin{tikzcd}[column sep=huge,row sep=large]
(\cte,\cto)
  \arrow[r,"\cte+X^k\cto"]
  \arrow[d,dashed,"\rep"']
&
\ct'_P
  \arrow[d,dashed,"\rep"]
\\
\bigl(\Pe(X^{2k}),\Po(X^{2k})\bigr)
  \arrow[r,"A+B\mapsto A+X^kB"']
&
P(X^k)
\end{tikzcd}
\]

\section{文档 2：大环密文到小环密文}

\subsection{问题是什么}

文档 2 的输入是一个大环密文
\[
\ct=(c_0(X),c_1(X))\in \Rbig^2
\]
满足
\[
 c_0(X)+s(X)c_1(X)=\Delta m(X^2)+e(X).
\]
这里 $m(X^2)$ 只含偶次幂。

目标是构造小环密文
\[
\ct'=(d_0(Y),d_1(Y))\in \Rsmall^2
\]
使得
\[
 d_0(Y)+s'(Y)d_1(Y)=\Delta m(Y)+e'(Y).
\]

也就是说，目标是把加密 $m(X^2)$ 的大环密文，转换成加密 $m(Y)$ 的小环密文。

\[
\begin{tikzcd}[column sep=huge,row sep=large]
\text{大环密文 } \ct=(c_0(X),c_1(X))
  \arrow[d,dashed,"\text{decrypt under }s(X)"']
  \arrow[r,"\Comp"]
&
\text{小环密文 } \ct'=(d_0(Y),d_1(Y))
  \arrow[d,dashed,"\text{decrypt under }s'(Y)"]
\\
\Delta m(X^2)+e(X)
  \arrow[r,"X^2\mapsto Y"']
&
\Delta m(Y)+e'(Y)
\end{tikzcd}
\]

\subsection{$\tau$ 映射是什么}

定义
\[
\tau:\Rbig\to \Rsmall.
\]
对任意
\[
F(X)=\sum_{j=0}^{2n-1}f_jX^j,
\]
令
\[
\tau(F)(Y)=\sum_{i=0}^{n-1}f_{2i}Y^i.
\]
也就是说，$\tau$ 做的是：取偶数下标系数，然后重新编号。

\[
\begin{tikzcd}[column sep=huge,row sep=large]
F(X)=f_0+f_1X+f_2X^2+f_3X^3+\cdots+f_{2n-1}X^{2n-1}
  \arrow[d,"\tau"']
\\
\tau(F)(Y)=f_0+f_2Y+f_4Y^2+\cdots+f_{2n-2}Y^{n-1}
\end{tikzcd}
\]

数组形式如下：

\[
\begin{tikzpicture}[>=Latex, node distance=1.1cm]
\node (top) {
\(\bigl[f_0,\ f_1,\ f_2,\ f_3,\ f_4,\ f_5,\ \cdots,\ f_{2n-2},\ f_{2n-1}\bigr]\)
};
\node[below=of top] (bottom) {
\(\bigl[f_0,\ f_2,\ f_4,\ \cdots,\ f_{2n-2}\bigr]\)
};
\draw[->] (top) -- node[right] {\(\tau\)：取偶数下标并重新编号} (bottom);
\end{tikzpicture}
\]

注意：文档 2 中 $\tau$ 作用在 $c_0,c_1$ 上时，并不是解密。$c_0,c_1$ 是密文分量，但它们本身是公开多项式数组，所以可以直接读系数并取偶数位置。

\subsection{公开拆分密文分量}

计算
\[
A=\tau(c_0),
\qquad
B=\tau(c_1),
\]
\[
C=\tau(X^{-1}c_0),
\qquad
D=\tau(X^{-1}c_1).
\]
因为在 $\Rbig=\ZZ_q[X]/(X^{2n}+1)$ 中
\[
X^{-1}=-X^{2n-1},
\]
所以乘以 $X^{-1}$ 只是公开的循环移位和变号。

图示如下：

\[
\begin{tikzcd}[column sep=huge,row sep=large]
c_0(X)
  \arrow[r,"\tau"]
  \arrow[d,"X^{-1}\cdot"']
&
A(Y)=\tau(c_0)
\\
X^{-1}c_0(X)
  \arrow[r,"\tau"']
&
C(Y)=\tau(X^{-1}c_0)
\end{tikzcd}
\qquad
\begin{tikzcd}[column sep=huge,row sep=large]
c_1(X)
  \arrow[r,"\tau"]
  \arrow[d,"X^{-1}\cdot"']
&
B(Y)=\tau(c_1)
\\
X^{-1}c_1(X)
  \arrow[r,"\tau"']
&
D(Y)=\tau(X^{-1}c_1)
\end{tikzcd}
\]

于是有恒等式
\[
c_0(X)=A(X^2)+XC(X^2),
\]
\[
c_1(X)=B(X^2)+XD(X^2).
\]
这里 $A,B$ 是偶部，$C,D$ 是奇部重新编号后的结果。

\subsection{同时拆分私钥}

由于解密关系里有 $s(X)c_1(X)$，所以私钥也必须拆成偶部和奇部：
\[
s(X)=s_0(X^2)+Xs_1(X^2).
\]

\[
\begin{tikzcd}[column sep=large,row sep=large]
c_0(X)
  \arrow[r,"\text{拆成偶部/奇部}"]
&
A(X^2)+XC(X^2)
\\
c_1(X)
  \arrow[r,"\text{拆成偶部/奇部}"]
&
B(X^2)+XD(X^2)
\\
s(X)
  \arrow[r,"\text{拆成偶部/奇部}"]
&
s_0(X^2)+Xs_1(X^2)
\end{tikzcd}
\]

这一步是文档 2 的关键：不是只压缩明文形式，还要把整个解密关系改写到小环中。

\subsection{重写解密方程}

原始解密式为
\[
c_0+s c_1=\Delta m(X^2)+e(X).
\]
代入
\[
c_0=A+XC,
\qquad
c_1=B+XD,
\qquad
s=s_0+Xs_1,
\]
这里为了简洁省略 $A(X^2),B(X^2)$ 等记号。

展开：
\[
\begin{aligned}
c_0+s c_1
&=A+XC+(s_0+Xs_1)(B+XD)\\
&=A+XC+s_0B+Xs_0D+Xs_1B+X^2s_1D.
\end{aligned}
\]
按偶部和奇部分组：
\[
\text{偶部：}\quad A+s_0B+X^2s_1D,
\]
\[
\text{奇部：}\quad X(C+s_0D+s_1B).
\]

图示如下：

\[
\begin{tikzcd}[column sep=huge,row sep=large]
c_0+s c_1
  \arrow[r,"\substack{c_0=A+XC\\ c_1=B+XD\\ s=s_0+Xs_1}"]
&
A+XC+(s_0+Xs_1)(B+XD)
  \arrow[d,"\text{展开}"]
\\
&
A+XC+s_0B+Xs_0D+Xs_1B+X^2s_1D
  \arrow[d,"\text{分离偶部和奇部}"]
\\
&
\begin{cases}
A+s_0B+X^2s_1D,\\
X(C+s_0D+s_1B).
\end{cases}
\end{tikzcd}
\]

由于 $\Delta m(X^2)$ 没有奇次项，若写
\[
e(X)=e_{\mathrm{even}}(X^2)+Xe_{\mathrm{odd}}(X^2),
\]
则比较偶部、奇部得到
\[
A+s_0B+Y s_1D=\Delta m(Y)+\nu_0,
\]
\[
C+s_0D+s_1B=\nu_1.
\]
其中 $\nu_0$ 来自原噪声的偶部，$\nu_1$ 来自原噪声的奇部。

\subsection{模块密文}

上面的两个方程可以写成矩阵形式：
\[
\begin{pmatrix}
A\\ C
\end{pmatrix}
+
\begin{pmatrix}
B & YD\\
D & B
\end{pmatrix}
\begin{pmatrix}
s_0\\ s_1
\end{pmatrix}
=
\begin{pmatrix}
\Delta m(Y)+\nu_0\\
\nu_1
\end{pmatrix}.
\]

图示如下：

\[
\begin{tikzcd}[column sep=huge,row sep=large]
\begin{pmatrix}
A\\
C
\end{pmatrix}
+
\begin{pmatrix}
B & YD\\
D & B
\end{pmatrix}
\begin{pmatrix}
s_0\\
s_1
\end{pmatrix}
  \arrow[r,"\text{decrypts to}"]
&
\begin{pmatrix}
\Delta m(Y)+\nu_0\\
\nu_1
\end{pmatrix}
\end{tikzcd}
\]

第一行
\[
A+s_0B+s_1(YD)=\Delta m(Y)+\nu_0
\]
携带目标消息。因此
\[
(t,u,v)=(A,B,YD)
\]
可以看成一个绑定向量私钥 $(s_0,s_1)$ 的三元素模块密文。

第二行
\[
C+s_0D+s_1B=\nu_1
\]
只对应近似零关系，通常不作为最终输出。

\subsection{正确的 key switching}

现在要把三元素模块密文
\[
(t,u,v)
\]
从向量私钥 $(s_0,s_1)$ 切换到普通小环私钥 $s'$。

也就是说，从
\[
t+s_0u+s_1v\approx \Delta m(Y)
\]
转换到
\[
d_0+s'd_1\approx \Delta m(Y).
\]

图示如下：

\[
\begin{tikzcd}[column sep=huge,row sep=large]
(t,u,v)=(A,B,YD)
  \arrow[d,dashed,"t+s_0u+s_1v"']
  \arrow[r,"\KS_{(s_0,s_1)\to s'}"]
&
(d_0,d_1)
  \arrow[d,dashed,"d_0+s'd_1"]
\\
\Delta m(Y)+\nu_0
  \arrow[r,"+\ \eta_{\mathrm{ks}}"']
&
\Delta m(Y)+\nu_0+\eta_{\mathrm{ks}}
\end{tikzcd}
\]

为了和解密约定
\[
d_0+s'd_1
\]
一致，切换密钥应满足
\[
b_{0,\ell}=G^\ell s_0-a_{0,\ell}s'+\epsilon_{0,\ell},
\]
\[
b_{1,\ell}=G^\ell s_1-a_{1,\ell}s'+\epsilon_{1,\ell}.
\]

把 $u,v$ 做 gadget 分解：
\[
u=\sum_{\ell=0}^{L-1}u_\ell G^\ell,
\qquad
v=\sum_{\ell=0}^{L-1}v_\ell G^\ell.
\]
初始化
\[
d_0=t,
\qquad
d_1=0.
\]
然后对每个 $\ell$ 更新
\[
d_0\leftarrow d_0+u_\ell b_{0,\ell}+v_\ell b_{1,\ell},
\]
\[
d_1\leftarrow d_1+u_\ell a_{0,\ell}+v_\ell a_{1,\ell}.
\]
最后输出
\[
\ct'=(d_0,d_1).
\]

正确性为
\[
\begin{aligned}
d_0+s'd_1
&=t+\sum_\ell u_\ell(b_{0,\ell}+s'a_{0,\ell})
     +\sum_\ell v_\ell(b_{1,\ell}+s'a_{1,\ell})\\
&=t+\sum_\ell u_\ell(G^\ell s_0+\epsilon_{0,\ell})
     +\sum_\ell v_\ell(G^\ell s_1+\epsilon_{1,\ell})\\
&=t+s_0u+s_1v+\eta_{\mathrm{ks}}.
\end{aligned}
\]
其中
\[
\eta_{\mathrm{ks}}
=
\sum_\ell u_\ell\epsilon_{0,\ell}
+
\sum_\ell v_\ell\epsilon_{1,\ell}.
\]

因此最终
\[
d_0+s'd_1=\Delta m(Y)+\nu_0+\eta_{\mathrm{ks}}.
\]

\subsection{噪声流向}

原始噪声分解为
\[
e(X)=e_{\mathrm{even}}(X^2)+Xe_{\mathrm{odd}}(X^2).
\]
输出噪声为
\[
e_{\mathrm{out}}=\nu_0+\eta_{\mathrm{ks}}.
\]
也就是说，原噪声的奇部不进入最终第一行输出；但是 key switching 会引入新噪声。

\[
\begin{tikzcd}[column sep=large,row sep=large]
\Delta m(X^2)+e(X)
  \arrow[r,"e=e_{\mathrm{even}}(X^2)+Xe_{\mathrm{odd}}(X^2)"]
&
\Delta m(X^2)+e_{\mathrm{even}}(X^2)+Xe_{\mathrm{odd}}(X^2)
  \arrow[d,"\text{投影到偶子环}"]
\\
&
\Delta m(Y)+\nu_0
  \arrow[d,"\text{key switching}"]
\\
&
\Delta m(Y)+\nu_0+\eta_{\mathrm{ks}}
\end{tikzcd}
\]

因此压缩操作不是免费消噪。它丢掉了与目标消息无关的奇部噪声，但又新增了 key switching 噪声。

\section{两个过程如何组合}

\subsection{一层递归}

把文档 1 和文档 2 连起来，一层递归可以画成：

\[
\begin{tikzcd}[column sep=large,row sep=large]
\ctP \text{ encrypts } P(X)
  \arrow[r,"\text{文档 1：Split}"]
&
\begin{matrix}
\cte \text{ encrypts } \Pe(X^2)\\
\cto \text{ encrypts } \Po(X^2)
\end{matrix}
  \arrow[r,"\text{文档 2：Compress each branch}"]
&
\begin{matrix}
\ct'_e \text{ encrypts } \Pe(Y)\\
\ct'_o \text{ encrypts } \Po(Y)
\end{matrix}
  \arrow[r,"\text{继续递归}"]
&
\cdots
\end{tikzcd}
\]

文档 1 完成
\[
P(X)\mapsto \Pe(X^2),\Po(X^2),
\]
文档 2 完成
\[
f(X^2)\mapsto f(Y).
\]
所以文档 2 可以接在文档 1 后面，让下一层递归真正发生在小环中。

\subsection{是否需要合并回去}

如果后续算法仍然在小环中进行，则不需要马上合并回去。

如果最终需要恢复父层大环密文，则需要先升环：
\[
\Pe(Y)\mapsto \Pe(X^2),
\qquad
\Po(Y)\mapsto \Po(X^2),
\]
再合并：
\[
P(X)=\Pe(X^2)+X\Po(X^2).
\]

图示如下：

\[
\begin{tikzcd}[column sep=huge,row sep=large]
\bigl(\ct'_e,\ct'_o\bigr)
  \arrow[r,"\text{升环 }Y\mapsto X^2"]
&
\bigl(\widehat{\ct}_e,\widehat{\ct}_o\bigr)
  \arrow[r,"\widehat{\ct}_e+X\widehat{\ct}_o"]
&
\ct_{\mathrm{merge}}
\\
\bigl(\Pe(Y),\Po(Y)\bigr)
  \arrow[r,"Y\mapsto X^2"']
&
\bigl(\Pe(X^2),\Po(X^2)\bigr)
  \arrow[r,"A+XB"']
&
P(X)
\end{tikzcd}
\]

注意：合并回去通常只能恢复语义等价的新密文，不会恢复原来的密文本身。压缩时已经经过 key switching，并且奇部噪声没有携带到主输出。

\section{层数消耗与复杂度}

\subsection{模数链层数消耗}

在典型 CKKS/RNS 实现中，层数消耗通常来自 rescale 或主动 modulus switching。自同构、加减法、key switching、单项式乘法、公开系数重排通常不消耗层数，但会增加噪声或计算成本。

\[
\begin{tikzcd}[column sep=large,row sep=large]
\text{输入 level } \ell
  \arrow[r,"\text{文档 1：Auto/Add/Sub/Mul by }1/2"]
&
\text{通常仍在 level } \ell
  \arrow[r,"\text{文档 2：}\tau+\text{KeySwitch}"]
&
\text{通常仍在 level } \ell
  \arrow[r,"\text{可选 ModDown/Rescale}"]
&
\text{level } \ell-1
\end{tikzcd}
\]

\begin{center}
\begin{tabular}{llll}
\toprule
模块 & 操作 & 默认层数消耗 & 说明 \\
\midrule
文档 1 & 自同构 $\Auto_{1+2n}$ & 0 & 增噪声，需要 automorphism key \\
文档 1 & 加法/减法 & 0 & 需同 level \\
文档 1 & 乘精确常数 $1/2$ & 0 & 若用 scale 1 plaintext \\
文档 1 & 乘单项式 $-X^{N-k}$ & 0 & 移位加变号 \\
文档 2 & $\tau(c_0),\tau(c_1)$ & 0 & 公开系数操作 \\
文档 2 & $X^{-1}c_0,X^{-1}c_1$ & 0 & 公开移位加变号 \\
文档 2 & $\tau(X^{-1}c_0),\tau(X^{-1}c_1)$ & 0 & 公开系数操作 \\
文档 2 & $YD$ & 0 & 小环单项式移位 \\
文档 2 & gadget decomposition & 0 & 系数分解 \\
文档 2 & key switching $(s_0,s_1)\to s'$ & 0 & 默认不掉层，但增噪声 \\
任意阶段 & rescale / mod-down & 1 或更多 & 主动降几层就消耗几层 \\
任意阶段 & ciphertext-ciphertext multiplication & 通常 1 & 当前主流程没有此操作 \\
\bottomrule
\end{tabular}
\end{center}

因此，该组合流程在理想实现口径下可以做到
\[
\boxed{0\text{ 层消耗}},
\]
真正的代价主要是 key switching 的噪声与计算量。

\subsection{复杂度概览}

设大环维度为 $N=2n$，小环维度为 $n$，gadget 分解长度为 $L$。

\begin{center}
\begin{tabular}{lll}
\toprule
步骤 & 复杂度 & 说明 \\
\midrule
提取 $A,B,C,D$ & $O(N)$ & 公开数组重排 \\
计算 $YD$ & $O(n)$ & 小环单项式移位 \\
gadget 分解 $B,YD$ & $O(Ln)$ & 对两个小环多项式分解 \\
key switching & $O(Ln\log n)$ & 主导成本，小环 NTT 乘法 \\
输出密文大小 & 从 $R_N^2$ 到 $R_n^2$ & 大约减半 \\
\bottomrule
\end{tabular}
\end{center}

所以文档 2 的压缩可以理解为：
\[
\boxed{\text{用一次小环 key switching 的代价，换来后续计算在半维度小环上进行。}}
\]
如果后续还有很多递归或同态计算，这通常是值得的；如果压缩后马上结束，则未必划算。

\section{完整主流程大图}

最后把两个文档的完整一层操作放在一张图中。

\[
\begin{tikzcd}[column sep=large,row sep=large]
\ctP \text{ in } \Rbig
  \arrow[r,"\Auto_{1+2n}"]
&
\ct_- \text{ in } \Rbig
\\
\ctP,\ct_-
  \arrow[r,"\frac{1}{2}(\ctP+\ct_-)"]
  \arrow[d,"\frac{1}{2}(\ctP-\ct_-)"']
&
\cte \text{ encrypts } \Pe(X^2)
  \arrow[r,"\Comp"]
&
\ct'_e \text{ in } \Rsmall \text{ encrypts } \Pe(Y)
\\
\ctoy \text{ encrypts } X\Po(X^2)
  \arrow[r,"\cdot X^{-1}"]
&
\cto \text{ encrypts } \Po(X^2)
  \arrow[r,"\Comp"]
&
\ct'_o \text{ in } \Rsmall \text{ encrypts } \Po(Y)
\end{tikzcd}
\]

这张图的读法是：
\begin{enumerate}
  \item 通过自同构得到 $P(-X)$ 的密文；
  \item 半和得到偶分支；
  \item 半差得到带 $X$ 因子的奇分支；
  \item 乘 $X^{-1}$ 得到归一化奇分支；
  \item 对两个形如 $f(X^2)$ 的分支分别执行文档 2 的压缩；
  \item 得到真正的小环密文 $f(Y)$，可以继续递归。
\end{enumerate}

\section*{最后总结}

文档 1 的核心是：
\[
\boxed{
P(X)=\Pe(X^2)+X\Po(X^2),
\qquad
P(-X)=\Pe(X^2)-X\Po(X^2).
}
\]
它通过自同构、半和、半差在密文上得到奇偶分支。

文档 2 的核心是：
\[
\boxed{
\text{若 } c_0+s c_1=\Delta m(X^2)+e,
\text{ 则可构造 } d_0+s'd_1=\Delta m(Y)+e'.
}
\]
它通过公开的 $\tau$ 拆密文分量、拆私钥、构造模块密文，再通过 key switching 得到普通小环密文。

两者合起来，就是：
\[
\boxed{
\text{先在密文上做 FFT 式拆分，再把每个半长度分支真正压到小环继续递归。}
}
\]

\end{document}
